% ============================================================================
% "Plausible Is Not Recovered" -- academic research poster (tikzposter, A0).
%
% Compile:  pdflatex poster.tex   (two passes)
% Zero-install option: upload this folder to a new Overleaf project.
% Images live in ./assets/ and are found via \graphicspath below.
% ============================================================================
\documentclass[14pt, a0paper, portrait, margin=6mm, innermargin=8mm,
               blockverticalspace=8mm, colspace=7mm, subcolspace=5mm]{tikzposter}

\usepackage[utf8]{inputenc}
\usepackage[T1]{fontenc}
\usepackage{amsmath, amssymb}
\usepackage{booktabs}
\usepackage{graphicx}
\usepackage{xcolor}
\graphicspath{{assets/}}

% --------------------------------------------------------------------------
% Clean journal-style look: white background, black section headings with a
% thin rule underneath, no frames, colour used only to carry meaning.
% --------------------------------------------------------------------------
\usetheme{Default}
\usetikzlibrary{calc}

\definecolor{posterblue}{HTML}{1F3A63}
\definecolor{postercyan}{HTML}{12706B}
\definecolor{posterred}{HTML}{B0413E}
\definecolor{posterink}{HTML}{111111}
\definecolor{posterrule}{HTML}{333333}

\definecolor{postertint}{HTML}{EEF2F7}   % pale blue block fill
\definecolor{posterwarm}{HTML}{FBF4E6}   % pale sand block fill
\definecolor{posterline}{HTML}{C9D3E0}   % column separator

\colorlet{backgroundcolor}{white}
\colorlet{framecolor}{white}
\colorlet{blocktitlefgcolor}{posterblue}
\colorlet{blocktitlebgcolor}{white}
\colorlet{blockbodyfgcolor}{posterink}
\colorlet{blockbodybgcolor}{white}
\colorlet{titlefgcolor}{posterink}
\colorlet{titlebgcolor}{white}

\tikzposterlatexaffectionproofoff

% --------------------------------------------------------------------------
% Block styles: a coloured rule under every heading, plus a tinted variant
% used to lift the blocks that carry the argument.
% --------------------------------------------------------------------------
\defineblockstyle{ruled}{
  titlewidthscale=1, bodywidthscale=1, titleleft, titleoffsetx=0pt,
  titleoffsety=0pt, bodyoffsetx=0pt, bodyoffsety=0pt, bodyverticalshift=0mm,
  roundedcorners=0, linewidth=0pt, titleinnersep=3mm, bodyinnersep=3mm
}{
  \ifBlockHasTitle
    \draw[color=posterblue, line width=2.2pt]
      (blocktitle.south west) -- (blocktitle.south east);
  \fi
}
\defineblockstyle{tinted}{
  titlewidthscale=1, bodywidthscale=1, titleleft, titleoffsetx=0pt,
  titleoffsety=0pt, bodyoffsetx=0pt, bodyoffsety=0pt, bodyverticalshift=0mm,
  roundedcorners=0, linewidth=0pt, titleinnersep=3mm, bodyinnersep=3mm
}{
  \ifBlockHasTitle
    \fill[blockbodybgcolor, rounded corners=6pt]
      ($(blockbody.south west) + (-4mm,-4mm)$)
      rectangle ($(blockbody.north east |- blocktitle.north) + (4mm,3mm)$);
    \draw[color=posterblue, line width=2.2pt]
      (blocktitle.south west) -- (blocktitle.south east);
  \else
    \fill[blockbodybgcolor, rounded corners=6pt]
      ($(blockbody.south west) + (-4mm,-4mm)$)
      rectangle ($(blockbody.north east) + (4mm,4mm)$);
  \fi
}
\useblockstyle{ruled}

% Thin vertical rules in the gutters, drawn on the poster background.
\newcommand{\gutterrule}[1]{%
  \coordinate (grb) at ($(bottomleft)!#1!(bottomleft -| topright)$);
  \coordinate (grt) at ($(topright -| bottomleft)!#1!(topright)$);
  \draw[color=posterline, line width=1.1pt] ($(grb)!0.055!(grt)$) -- ($(grb)!0.905!(grt)$);
}
\definebackgroundstyle{columnrules}{
  \fill[backgroundcolor] (bottomleft) rectangle (topright);
  \gutterrule{0.268}\gutterrule{0.700}
}
\usebackgroundstyle{columnrules}

% Wrap a block in a tinted panel: \tintblock{fill}{title}{body}
\newcommand{\tintblock}[3]{{\useblockstyle{tinted}\colorlet{blockbodybgcolor}{#1}\block{#2}{#3}}}

% Small coloured kicker used as an in-block label, e.g. \kicker{Why it matters}
\newcommand{\kicker}[1]{{\color{posterred}\bfseries\large #1}\par\vspace{0.2em}}
\newcommand{\subhead}[1]{{\color{posterblue}\bfseries #1}\par\vspace{0.15em}}
% Figure caption inside a block
\newcommand{\figcap}[1]{\vspace{0.25em}\par{\small #1}}

\settitle{
  \begin{minipage}[c]{0.21\linewidth}\ \end{minipage}%
  \begin{minipage}[c]{0.56\linewidth}
    \centering
    {\large\color{posterblue} A C A D E M I C \quad R E S E A R C H \quad P O S T E R}\\[0.55em]
    {\bfseries\Huge Plausible Is Not Recovered}\\[0.45em]
    {\LARGE Evidence-Grounded Reconstruction and Predictive Uncertainty\\[0.2em]
     for Partial Fingerprints: SOCOFing $\rightarrow$ NIST SD302}\\[0.55em]
    {\Large R\u{a}zvan-Cristian Alecse \quad\textbullet\quad University of Bucharest
     \quad\textbullet\quad [[Competition]] \quad\textbullet\quad [[Date]]}
  \end{minipage}%
  \begin{minipage}[c]{0.21\linewidth}
    \raggedleft
    \includegraphics[width=0.80\linewidth]{ub_logo.png}
  \end{minipage}
}

\begin{document}

\maketitle

\begin{columns}

% ==========================================================================
% LEFT COLUMN -- framing, formalism, protocol
% ==========================================================================
\column{0.26}

\block{What this project is}{
  \small
  Give a network half a fingerprint and it will happily draw the other half. The
  drawing usually looks right. The question this project got stuck on is whether
  there is any honest way to find out if it \emph{is} right, and what happens
  to your conclusions when the only available answer key is itself an
  approximation.
  \vspace{0.3em}

  Everything runs on public data (SOCOFing, NIST SD302). Nothing here identifies
  anyone or touches a sensor.
}

\tintblock{posterwarm}{Core distinction}{
  \small
  \kicker{Three different things}
  \textbf{Recovery} of the true ridge field, \textbf{structural
  plausibility}, and \textbf{calibrated uncertainty} are three different
  properties. A model can gain one while losing another, and most
  published-looking gains in this problem are gains in plausibility.
  A result is therefore only as strong as the ground truth behind it.
}

\block{Problem formulation}{
  \small
  A complete print $X \in [0,1]^{H \times W}$ is observed through a binary
  mask $M$:
  \begin{equation*}
    Y = M \odot X ,\qquad M \in \{0,1\}^{H \times W}
  \end{equation*}
  The estimation target is not one image but a conditional law
  \begin{equation*}
    p_\theta\!\left(X_{\text{missing}} \mid Y, M\right).
  \end{equation*}
  Under severe occlusion many completions are structurally plausible while
  only one matches the unknown original. Fidelity, plausibility and uncertainty
  therefore have to be measured as separate quantities, never as one score.
}

\tintblock{postertint}{What counts as ground truth}{
  \small
  The central design problem of this project: on real forensic latents,
  no target is simultaneously \emph{exact}, \emph{real} and \emph{dense}.
  \vspace{0.3em}

  \small
  \begin{tabular}{@{}llll@{}}
    \toprule
    Track & Target & Exact & Real latent \\
    \midrule
    \textbf{A} & SOCOFing, synthetic mask     & yes & no  \\
    \textbf{B} & $X_{\text{pseudo}}$ registered exemplar & no  & yes \\
    \textbf{C} & held-out \texttt{quality=1} pixels & yes & yes \\
    \textbf{E} & degraded clean exemplar      & yes & no  \\
    \bottomrule
  \end{tabular}
  \vspace{0.4em}

  \small
  Track B is the only dense target on real data, and it is
  \emph{not} ground truth. Track C is real but sparse; Track E is exact
  but simulated. \textbf{The protocol is therefore triangulation:} a claim
  counts only when independent tracks C and E agree, and Track B alone is
  never allowed to carry a conclusion.
}

\block{Experimental setup}{
  \small
  \begin{itemize}
    \item 8 observed fractions, $r = 0.10 \dots 0.80$.
    \item 8 mask families: rectangles, irregular, central, peripheral,
      disconnected fragments, stripes, severe partial.
    \item Exact reinjection of observed pixels, so the model is never credited for
      copying what it was given.
    \item $K \in \{5,10,20,50\}$ samples; single, mean-$K$ and best-of-$K$ reported
      separately, best-of-$K$ always flagged as an oracle quantity.
  \end{itemize}
  \vspace{0.15em}

  \small Difficulty is driven by \emph{where} the evidence is, not only how
  much survives: a ring of peripheral fragments and one central disc remove the
  same number of pixels and are not remotely the same problem.
}

\block{Ridge structure and metrics}{
  \small
  Ridge orientation is axial, i.e. $\pi$-periodic, so angular error uses
  the doubled-angle form
  \begin{equation*}
    d(\theta, \hat\theta) = 1 - \cos\!\left[2(\theta - \hat\theta)\right],
  \end{equation*}
  estimated only inside the fingerprint foreground.
  \vspace{0.3em}

  \subhead{Pixel} MAE, PSNR, local SSIM.
  \subhead{Structural} orientation error, ridge-frequency error, coherence, curvature,
  minutiae density.
  \subhead{Probabilistic} uncertainty--error $\rho$, interval coverage, CRPS.
}

\block{Statistical protocol}{
  \small
  \begin{itemize}
    \item Paired per-image comparisons on identical masks and seeds.
    \item Aggregated \textbf{per subject} (29 subjects), so many images of
      one finger cannot masquerade as independent evidence.
    \item Wilcoxon signed-rank with Holm--Bonferroni correction; Cohen
      $d_z$ reported alongside every $p$.
    \item Friedman omnibus before pairwise claims.
    \item The test split stayed sealed until every model, threshold and
      hyperparameter was frozen, then was evaluated \textbf{once}.
  \end{itemize}
}

\tintblock{posterwarm}{Controls that could have falsified us}{
  \small
  \begin{itemize}
    \item \textbf{Permuted identity.} Feed another finger's ridges through
      the same mask: MAE worsens by $0.034$, $d_z = -2.17$,
      $p = 3.7\times10^{-9}$. The model is using the evidence, not merely
      the mask shape.
    \item \textbf{Non-rigid registration.} Thin-plate spline vs.\ affine
      improves leave-one-out residual by only $0.005$--$0.012$\,mm
      ($p = 0.07 / 0.14$): affine is not the bottleneck, so the pseudo-target
      problem is not a registration problem.
    \item \textbf{Data-consistency ablation.} Null by construction, and
      reported as null.
    \item \textbf{A bug we found and disclosed.} Brownian Bridge sampling
      divided by $(1-m_t)$, which reaches $10^{-4}$ at the first step and
      amplified error up to $10^{4}\times$, collapsing outputs. Rescaled
      ($s: 1.0 \rightarrow 0.3$, $\epsilon: 10^{-4} \rightarrow 0.05$),
      retrained and verified visually before any number from that model was
      reported.
  \end{itemize}
}

\block{Interpretation and next step}{
  \small
  \subhead{Positive evidence} diffusion restoration improves fidelity over
  deterministic baselines; structural losses preserve ridge flow; multiple
  samples expose genuine ambiguity; boundary continuity gives a real,
  large, reproducible gain. \\[0.3em]
  \subhead{Negative evidence} uncertainty is miscalibrated for
  latent-variable models; fine-tuning overfits registration artefacts;
  latent diffusion underperforms at this resolution; ensembling did not
  beat the best single model. \\[0.4em]
  \subhead{Next} progressive structure-conditioned residual diffusion at
  native ridge resolution, with elastic-registration uncertainty and
  subject-level calibration folded into the loss.
}

\block{Is it ridges, or is it texture?}{
  \small
  A reconstruction can be locally convincing and topologically empty.
  Comparing reconstruction against real ridges inside the \emph{same}
  region, over 284 images, with 95\% CIs excluding zero on all three:
  \vspace{0.3em}

  \centering
  \begin{tabular}{@{}lrl@{}}
    \toprule
    Reconstruction $-$ real & Value & Reading \\
    \midrule
    Orientation coherence & $+0.181$ & too uniform \\
    Orientation curvature & $-0.022$ & too straight \\
    Minutiae density      & $-9.78$ / 1000\,px & about half \\
    \bottomrule
  \end{tabular}
  \vspace{0.4em}

  \raggedright
  The model paints a smooth, well-oriented ridge field and drops the endings and
  bifurcations that carry identity. That is the quantitative form of the ``false
  ridges'' failure, and it is invisible to MAE and SSIM, which improve while it happens.
}

\block{Where the evidence stops}{
  \small
  The SD302 target is another impression of the same finger, so even the honest
  supervision here is structural, never pixel-exact. The held-out
  \texttt{quality=1} pixels are faint \emph{existing} traces rather than truly
  absent regions, which means the hardest version of this problem is still
  open: \textbf{recovery of ridges with no observable trace has never been
  measured on a real latent}, by this project or anyone else. And 29 validation
  subjects with 217 test images carry paired subject-level tests, not claims
  about a population.
}

\block{Future directions}{
  \small
  \begin{itemize}
    \item \textbf{Does the reversal change the model \emph{ranking}?} Shown here
      for one checkpoint pair. Ranking all eleven models by $X_{\text{pseudo}}$ and
      by real pixels, then correlating the two orderings, turns a caution into a
      number: how often does the convenient target pick the wrong model?
    \item \textbf{Minutiae matching, not density.} Counting shows detail is
      missing; matching would say whether what is synthesised is forensically
      meaningful.
    \item \textbf{Native ridge resolution.} Simply raising the working
      resolution to 256px made every structural metric worse, so the sampling
      argument now rests on native-resolution \emph{patches} alone; a full
      patch pipeline is the remaining way to settle it.
    \item \textbf{Registration uncertainty in the loss}, so the model stops being
      penalised for the target's own misalignment.
  \end{itemize}
}

\block{}{
  \fboxsep=10pt
  \colorbox{posterblue}{\parbox{0.955\linewidth}{
    \color{white}
    {\bfseries\Large Findings}\par\vspace{0.35em}
    \small
    \begin{itemize}
      \item On real held-out pixels, zero-shot transfer \textbf{beat} SD302
        fine-tuning on every metric, the opposite of the pseudo-target
        conclusion, and it held on the sealed test set.
      \item Across all eleven models, \textbf{no model won fidelity and uncertainty
        calibration together}: RePaint took MAE ($0.1008$), DDIM took calibration
        ($\rho = 0.276$) from second-worst MAE.
      \item The spatial CVAE's uncertainty \textbf{inverts sign} ($\rho = -0.153$):
        most confident where it is most wrong.
      \item A boundary phase-continuity loss, a single added term, gave one of the
        largest single-change gains measured ($d_z$ up to $2.66$, $p = 1.9\times10^{-8}$).
      \item Reconstructions carry roughly \textbf{half} the minutiae density of the
        ridges they replace: plausible texture, missing identity.
    \end{itemize}
  }}
}

\tintblock{posterwarm}{Can the damage be undone?}{
  \small
  If the harm comes from the target, fixing the target should fix the harm.
  Two ways to try, both evaluated on real held-out pixels.
  \vspace{0.3em}

  \subhead{Reweight it} Measure, per pixel, whether the latent and the
  registered exemplar agree on ridge direction, and weight the loss by that.
  Agreement averages only $0.586$ and ranges $0.21$--$0.87$, so the target
  really is unreliable and unevenly so. Orientation error $0.517
  \rightarrow 0.454$ ($d_z = 1.38$, $p = 4.5\times10^{-8}$):
  \textbf{37\% of the damage recovered.}

  \subhead{Replace it} Train instead on an \emph{exact} target built inside the
  SD302 domain, by degrading a clean exemplar and masking it, so nothing
  approximate enters the loss. Orientation $0.517 \rightarrow 0.402$, MAE
  $0.144 \rightarrow 0.113$, SSIM $0.205 \rightarrow 0.347$ ($d_z$ up to
  $2.87$): \textbf{two thirds recovered}, roughly double what reweighting
  buys.
  \vspace{0.25em}

  \textbf{Neither beats zero-shot.} A third of the gap survives an exact target
  in the right image domain, so ``the target is wrong'' is not the whole
  story. Either the simulated degradation does not resemble real latent
  degradation closely enough, or adapting to this domain is simply harmful at
  this data volume. Separating those needs a genuine degraded-and-clean pair,
  which SD302 does not provide.
}

% ==========================================================================
% CENTRE COLUMN -- the evidence
% ==========================================================================
\column{0.44}

\block{Track A: does the model reconstruct, or does it blur?}{
  \centering
  \includegraphics[width=1.0\linewidth]{grid_socofing.png}
  \figcap{Clean data, exact ground truth, four mask geometries. Structural
  losses (orientation, ridge-band, Gabor spectrum) buy large gains in
  ridge geometry at flat or slightly worse MAE. On disconnected
  fragments, orientation error falls $0.346 \rightarrow 0.211$ while MAE
  moves by $0.006$. \textbf{The best-looking reconstruction is not the
  lowest-MAE reconstruction}, which is the first warning that a single
  pixel metric cannot rank these models.}

  \vspace{0.3em}
  \raggedright\small
  Against the non-learned floor (900-image validation): gated convolution beats
  U-Net on every metric, paired and Holm-significant, but by small margins
  ($d_z\,0.13$--$0.31$). Structural losses on top improve all four at once
  (SSIM $d_z = 0.76$, $p = 2.8\times10^{-90}$).
}

\block{What a general-purpose image model does with this task}{
  \centering
  \includegraphics[width=1.0\linewidth]{generalist_baseline.png}
  \raggedright
  \small
  \vspace{0.3em}
  The same masked print and the same instruction were given to a
  general-purpose image generator, five times. It was told to preserve the
  visible pixels and not change them. It did not: it re-rendered the whole
  image, and the hidden region came back different every time.
  \vspace{0.3em}

  \textbf{What it produced looks like an excellent fingerprint.} Minutiae
  density $39.9$--$59.2$ per 1000\,px against $44.5$ in the truth; orientation
  coherence $0.55$--$0.66$ against $0.554$. On every "does this look real"
  statistic, it passes.
  \vspace{0.2em}

  \textbf{It is the wrong fingerprint.} Mean orientation error $0.427$, against
  $0.094$ for this project's constrained model on the same prints. Two
  completions score \emph{negative} SSIM. And the truth is an arch with no
  core, while the completions invent a whorl core: not a ridge error, a
  different pattern class, which forensically is a different finger.
  \vspace{0.2em}

  \footnotesize\emph{Not a like-for-like benchmark:} this project's models are
  constrained to the observed pixels by construction and the generalist model
  ignored that constraint, so the comparison measures task compliance, not
  architecture. Five completions of one print: an illustration, not a statistic.
}

\block{Why real forensic data has no exact ground truth}{
  \centering
  \includegraphics[width=1.0\linewidth]{registration.png}
  \figcap{The SD302 ``target'' is a \emph{different impression of the same
  finger}, aligned to the latent through 13 examiner-verified minutiae
  correspondences (RMSE $0.102$\,mm). Registration recovers pose, not
  ridge phase: $X_{\text{pseudo}}$ is structurally right and pixel-wise
  wrong. Training against it is training against a plausible lie.}

  \vspace{0.5em}
  \includegraphics[width=1.0\linewidth]{nist_pipeline.png}
  \figcap{The conditional pipeline on one real latent: observed
  quality-gated pixels $Y$, the registered pseudo-target, the
  reconstruction, the predicted visible support $P(\text{support}\mid Y,M)$
  that suppresses ridge hallucination into scanner background, the
  geometric confidence map, and the residual against
  $X_{\text{pseudo}}$, which is concentrated exactly where the two
  impressions genuinely differ.}
}

\tintblock{postertint}{\textcolor{posterred}{The central finding: the target can reverse the conclusion}}{
  \small
  Scored against $X_{\text{pseudo}}$, SD302 fine-tuning \textbf{improves} MAE and
  orientation, which is the result one would normally report. Scored against real held-out
  \texttt{quality=1} pixels, never touched by any loss, the \emph{same two checkpoints}
  reverse completely: worse on all four metrics (MAE $+0.040$, orientation $+0.058$,
  ridge-frequency $+0.202$, SSIM $+0.222$; all Holm $p < 10^{-4}$). A 6-model Friedman
  omnibus confirms the differences are real ($\chi^2 = 116.27$, $p = 1.9\times10^{-23}$),
  and the subject-mean ranking places the checkpoint that never saw SD302 \textbf{first}.
  \vspace{0.4em}

  \centering
  \includegraphics[width=0.98\linewidth]{chart_critical_finding.png}
  \figcap{Sealed test split, 217 images, evaluated once.}

  \vspace{0.3em}
  \raggedright
  \small
  \subhead{Triangulated three independent ways}
  (i) real held-out pixels; (ii) Track E exact synthetic ground truth, agreeing with one of
  the largest effect sizes in the project (SSIM $d_z = -2.78$, $p = 1.5\times10^{-8}$);
  (iii) the sealed test split ($0.105 \rightarrow 0.142$). \textbf{Consequence:} a measured
  improvement can be a property of the evaluation target rather than of the model.
  Fine-tuning on a registered different-impression target teaches the network to reproduce
  \emph{that} impression's texture: locally convincing, and wrong.
}

\tintblock{posterwarm}{Why zero-shot wins: it learns the finger, not the impression}{
  \small
  The reversal above says \emph{that} the target corrupts the result. This says
  \emph{why}, and it is the question the finding had left open.
  \vspace{0.3em}

  Pixels where the latent and the registered exemplar \emph{agree} cannot
  discriminate: following either one looks identical. So the test isolates
  pixels where the two impressions \textbf{conflict}, and asks whether the
  model's ridge orientation lands closer to the real latent or to
  $X_{\text{pseudo}}$. Zero-shot never saw $X_{\text{pseudo}}$, so its rate is
  what not-following looks like.
  \vspace{0.3em}

  \centering
  \begin{tabular}{@{}lrrr@{}}
    \toprule
    Held-out real pixels & Zero-shot & Fine-tuned & $d_z$ \\
    \midrule
    Follows $X_{\text{pseudo}}$ where impressions conflict & 0.243 & \textbf{0.432} & $-1.88$ \\
    Orientation error, conflicting pixels & 0.497 & \textbf{0.780} & $-1.75$ \\
    Orientation error, agreeing pixels & 0.342 & 0.390 & $-0.48$ \\
    \bottomrule
  \end{tabular}
  \vspace{0.35em}

  \raggedright
  Fine-tuning does not degrade reconstruction diffusely. It teaches the network
  to reproduce \emph{that impression's} ridge placement, and the cost lands
  precisely where that impression differs from the one being reconstructed: a
  78\% relative rise in following the target, error 57\% higher in exactly
  those pixels, and a much smaller gap everywhere else. All at
  $p < 4\times10^{-9}$, aggregated per subject.
}





% ==========================================================================
% RIGHT COLUMN -- models, uncertainty, controls, outlook
% ==========================================================================
\column{0.30}

\tintblock{posterwarm}{What this project contributes}{
  \small
  Not a better reconstructor. A harder test of whether any reconstructor is
  working.
  \vspace{0.35em}

  \begin{itemize}
    \item \textbf{The evaluation target can invert the conclusion.} Fine-tuning
      is a win against the registered exemplar and a loss against real pixels:
      MAE $0.105 \rightarrow 0.142$ on the sealed test set, worse on all four
      metrics, Friedman $\chi^2 = 116.27$ across six models. Checked three
      independent ways before it was believed.
    \item \textbf{Track E.} Exact ground truth manufactured inside SD302 by
      degrading a clean exemplar, so that claim can be tested without leaving
      the dataset or trusting the pseudo-target.
    \item \textbf{Uncertainty measured against distance from evidence}, not as
      one global number. DDIM decays gracefully ($0.527 \rightarrow 0.167$);
      the CVAE crosses zero ($+0.011 \rightarrow -0.297$) and grows confident
      exactly where it extrapolates furthest.
    \item \textbf{Three numbers that catch a convincing fake.} Reconstructions
      run $+0.18$ orientation coherence, $-0.02$ curvature and \emph{half} the
      minutiae density of the ridges they replace, while MAE and SSIM improve.
    \item \textbf{A permuted-identity control.} Another finger's ridges through
      the same mask and the model degrades ($d_z = -2.17$), so the skill is not
      mask-shape bookkeeping.
  \end{itemize}
  The models are standard. The protocol is not.
}

\block{Datasets and leakage control}{
  \small
  \subhead{SOCOFing} 6{,}000 images, 600 subjects, clean, synthetic masks. \\
  \subhead{NIST SD302} real forensic latents (A/B/D), 200 subjects,
  official EFS quality maps and minutiae correspondences. \\
  \subhead{Sealed test} 217 images from subjects unseen during any
  training, tuning or model selection.
  \vspace{0.35em}

  Splits are subject- and finger-disjoint; thresholds, Gabor bank and confidence
  weights were calibrated on validation only; nothing was re-selected after the test run.
}

\block{Model family}{
  \small
  \subhead{Deterministic} nearest-observed interpolation (non-learned
  floor), U-Net, gated convolution, $+$ structural losses. \\
  \subhead{Latent-variable} global and spatial CVAE. \\
  \subhead{Generative restoration} conditional DDPM, DDIM, RePaint,
  residual DDPM, latent diffusion (plain and structure-guided), conditional
  flow matching, Brownian Bridge Diffusion.
  \vspace{0.35em}

  All conditioning is \textbf{ControlNet-style}: the auxiliary signal is injected at
  every encoder scale through zero-initialised $1\times1$ projections, not concatenated
  once at the input.
}

\block{Does uncertainty stay honest?}{
  \centering
  \includegraphics[width=0.98\linewidth]{chart_uncertainty_distance.png}
  \raggedright
  \small
  \vspace{0.2em}
  Predictive spread is only useful if it grows where the model is
  actually wrong. Stratified by distance to the nearest examiner-verified
  correspondence (validation split), DDIM leads at every band and degrades
  gracefully; RePaint decays to noise; \textbf{the CVAE changes sign},
  becoming systematically overconfident exactly where it extrapolates
  furthest. For anything examiner-facing, that failure mode is
  disqualifying regardless of mean-image quality.
  \vspace{0.3em}

  The sealed test set reproduces the ordering without stratification:
  DDIM $+0.276$, RePaint $+0.135$, residual DDPM $+0.037$, flow matching
  $+0.021$, CVAE $\mathbf{-0.153}$. Uncertainty quality and fidelity rank
  the family in almost opposite orders.
}

\block{Full results matrix: sealed test set}{
  \small
  All eleven models below were run on the \emph{same} 217-image sealed
  test split, once, after every design choice was frozen.
  \vspace{0.35em}

  \small
  \begin{tabular}{@{}lccccc@{}}
    \toprule
    Model & $K$ & MAE $\downarrow$ & SSIM $\uparrow$ & Orient.\ $\downarrow$ & $\rho_{\sigma,e} \uparrow$ \\
    \midrule
    \multicolumn{6}{@{}l}{\itshape Deterministic} \\
    Zero-shot transfer         & n/a  & 0.1054 & \textbf{0.4884} & \textbf{0.2764} & n/a \\
    Fine-tuned (supp.+spec.)   & n/a  & 0.1416 & 0.2802 & 0.3237 & n/a \\
    $+$ boundary continuity    & n/a  & 0.1346 & 0.3036 & 0.3090 & n/a \\
    \midrule
    \multicolumn{6}{@{}l}{\itshape Probabilistic (predictive mean)} \\
    \textbf{RePaint}           & 5    & \textbf{0.1008} & 0.4382 & 0.2838 & 0.135 \\
    Residual DDPM              & 10   & 0.1126 & 0.3916 & 0.2856 & 0.037 \\
    Learned ensemble           & 5    & 0.1151 & 0.3966 & n/a & n/a \\
    CVAE, spatial              & 50   & 0.1378 & 0.3647 & 0.2952 & $-0.153$ \\
    DDIM-20                    & 10   & 0.1453 & 0.3013 & 0.2932 & \textbf{0.276} \\
    Flow matching, 8 steps     & 10   & 0.1545 & 0.3157 & 0.2966 & 0.021 \\
    Latent DDPM                & 10   & 0.2111 & 0.2411 & 0.3074 & n/a \\
    Latent DDPM, guided        & 10   & 0.2143 & 0.2357 & 0.3077 & n/a \\
    \bottomrule
  \end{tabular}
  \vspace{0.4em}

  \small
  \textbf{No row dominates.} RePaint takes fidelity at $35$--$60\times$ DDIM's
  per-sample cost; DDIM takes calibration by a wide margin ($0.276$ at $K{=}10$, $0.338$
  at $K{=}50$) from near the bottom on MAE; the zero-shot checkpoint still holds the best
  SSIM and orientation error. Flow matching reaches DDIM-comparable ridge-frequency error
  ($0.499$ vs.\ $0.503$) using \textbf{8 sampling steps instead of 500}. Latent diffusion
  is a clean negative result in both variants. The learned ensemble beats simple averaging
  ($0.1151$ vs.\ $0.1198$) but loses to its own best member.
  \vspace{0.3em}

  \centering
  \includegraphics[width=0.98\linewidth]{chart_model_comparison.png}
}

\tintblock{postertint}{A boundary phase-continuity loss}{
  \small
  Penalise the discontinuity of the ridge signal exactly where the model
  hands off from observed evidence to synthesis, measured along the local
  ridge tangent $\tau(p)$ from the target's structure tensor (term proposed in
  external review; implemented and measured here):
  \begin{equation*}
    \mathcal{L}_{\partial} =
    \frac{1}{|\partial M|}\sum_{p \in \partial M}
    \left|\, \nabla_{\tau(p)}\hat X_p - \nabla_{\tau(p)} X_p \,\right|
  \end{equation*}
  Added to the full recipe, identical weights otherwise:
  MAE $0.148 \rightarrow 0.141$, SSIM $0.281 \rightarrow 0.301$,
  orientation $0.337 \rightarrow 0.324$; $d_z = 1.01$--$2.66$, Holm $p$
  down to $1.9\times10^{-8}$ (287 images, 29 subjects). Three of four
  metrics improve with large effect sizes and the fourth moves the same
  way. It attacks precisely the seam the ``false ridges'' diagnostic
  identified.
}





\tintblock{postertint}{Knowing when not to trust it}{
  \centering
  \includegraphics[width=0.99\linewidth]{chart_reliability.png}
  \raggedright
  \small
  \vspace{0.2em}
  Three results, all aimed at the same question: given a reconstruction and no
  target, can we say which parts of it to believe?
  \vspace{0.3em}

  \subhead{A second opinion predicts the model's own error}
  The disagreement between the model's ridge orientation and a harmonic
  completion of the observed field predicts that model's true error at
  Spearman $\rho = 0.67$ per image and $0.71$ per subject, against $0.276$ for
  the best predictive-uncertainty signal in the whole study. It is not
  circular: controlling for the geometric solver's own accuracy leaves
  $\rho = 0.635$ ($p = 1.9\times10^{-33}$).

  \subhead{Abstention converts that into error reduction}
  Answering for only 60\% of the missing region cuts orientation error by
  \textbf{30\%}. Random abstention is flat, so the gain is selection, not
  arithmetic. Cross-validated over whole subjects; the oracle gap is large and
  reported rather than hidden.

  \subhead{And it can be made a guarantee}
  Split conformal over disjoint subject folds gives distribution-free bounds
  whose empirical coverage matches the target to the third decimal and holds
  per subject ($0.837$--$0.970$ at a $90\%$ target). \textbf{The limit is the
  reconstruction, not the method:} below $30^{\circ}$ tolerance nothing can be
  certified at all; at $45^{\circ}$, $41\%$ of regions can, with half the error
  of the rest.
}








\end{columns}

\block{}{
  \small
  \textbf{Safety and interpretation.} A plausible generated ridge field must never be
  described as recovery of an individual's true missing fingerprint without evidence. This
  project is limited to academic benchmark reconstruction and uncertainty analysis, and its
  central result is a caution about evaluation, not a capability claim.
  \hfill
  \textbf{Contact:} [[email]] \quad \textbf{Project:} [[URL / QR code]]
  \vspace{0.35em}

  \centering \footnotesize
  [1] Ho et al., \emph{Denoising Diffusion Probabilistic Models}, NeurIPS 2020 \quad\textbullet\quad
  [2] Lugmayr et al., \emph{RePaint}, CVPR 2022 \quad\textbullet\quad
  [3] Rombach et al., \emph{Latent Diffusion}, CVPR 2022 \quad\textbullet\quad
  [4] Li et al., \emph{BBDM}, CVPR 2023 \quad\textbullet\quad
  [5] NIST Special Database 302 \quad\textbullet\quad
  [6] Lipman et al., \emph{Flow Matching for Generative Modeling}, ICLR 2023
}

\end{document}
